\section{问题重述与分析}

\subsection{问题背景}

某山区乡村拥有平旱地、梯田、山坡地、水浇地、普通大棚和智慧大棚等多类耕地，不同地块在面积、可种植作物、种植季次和生产条件上存在明显差异。当地可选择的农作物包括粮食作物、蔬菜和食用菌等，不同作物的亩产量、种植成本、销售价格和适宜耕地类型也不相同。由于耕地资源有限，且作物之间存在轮作、连作、季节安排和销售消化能力等约束，简单按照单作物收益高低安排种植，容易造成土地利用冲突、产量过剩或农作制度不合理。

题目要求在 2024--2030 年规划期内，为该乡村制定合理的农作物种植方案。模型既要满足地块类型与作物适宜性的基本规则，又要考虑豆类轮作、同地块连作限制、每季种植面积约束等实际种植要求；同时，还需处理未来销售量、价格、成本和亩产等因素的不确定变化，并进一步分析不同作物之间可能存在的替代性和互补性关系。因此，本题本质上是一个带有多时期、多地块、多作物和多季次约束的农业种植优化问题。

\subsection{问题重述}

根据题目给出的地块信息、作物信息、2023 年种植情况和统计数据，需要解决以下三个层次的问题。

第一，假定 2024--2030 年各作物的预期销售量、亩产量、种植成本和销售价格均保持与 2023 年相同，需要在满足农作制度约束的前提下，为未来七年制定种植计划。对于超过预期销售量的产量，题目给出两种处理方式：一种是超出部分滞销，不能产生销售收入；另一种是超出部分按正常价格的 50\% 降价出售。需要分别在两种情形下确定各地块、各作物、各季次、各年份的种植面积，并填写相应的结果表。

第二，在未来七年中，市场需求、种植成本、销售价格和亩产量均可能发生变化。其中，小麦和玉米需求呈增长趋势，其他作物需求存在波动；种植成本整体上升；蔬菜价格有增长趋势，食用菌价格存在下降趋势；亩产量还会受到自然波动和露天耕地灾害影响。需要在这些不确定因素下，重新制定 2024--2030 年的种植方案，使方案不仅具有较高期望收益，还能兼顾风险控制。

第三，在进一步考虑作物之间关系的情况下，需要分析不同农作物在销售量、价格和成本等方面的统计相关性，并据此刻画作物之间的替代关系和互补关系。替代关系会影响作物的需求分配，互补关系则可能影响后茬作物的种植成本。需要在问题二随机规划框架的基础上引入这些相关性机制，形成更贴近实际农业经营的种植优化模型，并比较其与前两问方案的差异。

\subsection{问题分析}

本题的求解需要先将原始附件整理为可用于优化模型的数据结构。附件1给出了地块类型、地块面积和作物适宜种植规则，可据此建立“地块--作物--季次”的可行种植集合；附件2给出了2023年的实际种植情况和统计数据，可用于生成基准销售量、亩产量、种植成本和销售价格等参数；附件3为结果填报模板，需要在模型求解后按年份、季次和地块回填各作物种植面积。因此，建模前应先完成数据清洗与参数构造，包括合并单元格信息补全、价格区间取均值、说明行与空行剔除、2023年种植锚点提取，以及智慧大棚第一季缺失参数的合理平补。

在完成数据预处理后，问题一可作为全题的确定性基准模型。首先，根据地块类型和作物适宜规则构造可行种植集合，排除不符合种植制度的地块--作物--季次组合；然后，以七年总利润最大为目标，分别刻画超产滞销和超产半价销售两种收入口径；最后，引入种植面积变量和0-1种植指示变量，将面积限制、最小种植比例、连作禁止、豆类轮作和作物集中度等要求统一写入混合整数线性规划模型。该模型的作用是给出在基准年份参数不变条件下的种植方案，并为后续随机模型提供可行的种植拓扑基础。

问题二在问题一的基础上进一步考虑未来不确定性。首先，依据题目给出的增长和波动规则，分别生成销售量、种植成本、销售价格、亩产波动和露天灾害等随机情景；然后，将种植面积和种植拓扑视为规划期开始前确定的第一阶段决策，将各情景下的正常销售量和折价销售量作为第二阶段决策；最后，引入CVaR风险度量，在期望利润最大化的同时控制尾部不利情景的影响。这样可以避免模型只追求平均收益，而忽略灾害或价格波动下的经营风险。

问题三继续放宽“作物相互独立”的假设，将作物之间的统计关系纳入种植决策。首先，利用作物需求、价格、成本和波动特征进行标准化处理，并通过Spearman相关系数和K-means聚类识别作物之间的替代性和功能分组；然后，对同簇且负相关程度较强的作物引入需求转移机制，刻画一种作物供给不足时替代作物可能承接部分需求的情形；最后，利用豆类作物的前茬效应构造互补性成本折减机制，将其嵌入问题二的随机规划框架。由此得到的扩展模型能够同时反映市场替代关系和农艺互补关系，使种植方案更贴近实际农业经营。

综上，本文的建模路线为：先通过数据预处理建立统一参数体系，再构造确定性混合整数线性规划作为基准方案，随后引入随机情景和CVaR形成风险约束下的两阶段随机规划，最后加入作物替代与互补机制完成扩展优化。三个问题由确定性到随机性、由单作物独立到作物关联逐层递进，便于比较不同现实因素对种植策略和收益结构的影响。
